\section{The original-data variational identity}
\label{var:section}

We now assume the connected-horizon hypotheses of
Definition~\ref{def:rigidity-class} and equality in
Theorem~\ref{thm:numerical}, with $m>|Q|$.  All objects in this section
belong to the original exterior $(\Omega,g,K,E)$.  In particular, no
equality statement about one of the auxiliary metrics is used.  We first
differentiate the full enclosing-area infimum, then separate the active
constraint inequalities from a direction that would improve equality.
The resulting multipliers are initially measures.  Their regularity and
the elimination of their possible singular area sheets are separate
steps.
Decrease the asymptotic exponent, if necessary, so that
$(n-2)/2<q<n-2$; this retains every original hypothesis.

Write $A=\Area_g(S)$, $k=n-1$, and
\begin{equation}
 r_h=(A/\omega)^{1/k},\qquad
 b^0=\frac{\mathcal E_{\mathrm{ADM}}}{m},\qquad
 b^i=-\frac{(\mathbf P_{\mathrm{ADM}})_i}{m},\qquad
 c=\frac{k-1}{r_h}\left(1-\frac{Q^2}{r_h^{2(k-1)}}\right)>0.
 \label{var:constants}
\end{equation}
The vector $b$ is fixed when taking variations.  Set
$\mathcal B=b^0\mathcal E_{\mathrm{ADM}}+\sum_i b^i
(\mathbf P_{\mathrm{ADM}})_i$.  Thus $m'=\mathcal B'$ at the original
data.  On the outer, nonextremal area branch, differentiation at fixed
$Q$ gives
\begin{equation}
 \frac{\dd}{\dd A}\frac12
 \left(r_h^{k-1}+Q^2r_h^{-(k-1)}\right)
 =\frac{c}{2k\omega}.
 \label{var:area-coefficient}
\end{equation}
The strictly positive coefficient in this formula will be essential.

Throughout the variations, an electric field is represented by the
closed flux form $\alpha=i_E\vol_g$.  We also allow a closed, compactly
supported two-form $D_m$, initially zero, representing the spatial
Faraday form divided by $c_n$.  For these extended data put
\begin{equation}
 \begin{split}
 \mu_m&=\frac12(R_g+\tau^2-|K|^2)
              -c_n^2(|E|^2+|D_m|^2),\\
 J_m&=\operatorname{div}(K-\tau g)-2c_n^2D_m(\,\cdot\,,E).
 \end{split}
 \label{var:magnetic-constraints}
\end{equation}
Only the strictly feasible numerical extension of
Corollary~\ref{flux:magnetic-extension} will be applied to data with
$D_m\ne0$.  No magnetic equality theorem is assumed.  Compactly
supported closed variations of either form, including tests supported
up to $S$, do not change $Q$, by Stokes' theorem on a truncation.

\subsection{The full area envelope}

Choose a compact smooth filling $O$ behind $S$, with a smooth extension
of the metric, and denote the resulting boundary-free manifold by
$\mathcal M$.  The filling is used only to define perimeter
competitors; no constraint inequality is required there.  A filled
competitor is a bounded finite-perimeter set $F\subset\mathcal M$
containing $O$ up to null sets.  Its perimeter counts the part coinciding
with $\partial O=S$.

\begin{proposition}[Area envelope and regular patches]
\label{var:area-envelope}
The least perimeter of filled competitors equals the full smooth-cut
infimum $a$.  The family $\mathcal F_*$ of minimizing filled sets,
considered modulo null sets, is compact in local $L^1$, and all its
members are supported in a common compact set.  Their exterior end
components have the full-boundary convention of
Definition~\ref{def:full-cut}.

Each minimizing frontier is smooth and minimal at its free regular
points.  Its possible singular set has Hausdorff dimension at most
$n-8$.  All contact points with $S$ are regular of class $C^{1,\gamma}$
for some $\gamma>0$, and their graph representations belong to
$W^{2,p}$ for every finite $p$.  At an interior regular point of one
frontier, any other minimizing frontier through the point coincides
with it locally.  On a smaller neighborhood the minimizing frontiers
that meet that neighborhood are disjoint single graphs, with uniform
smooth interior estimates.

For a differentiable metric path with $g'_0=h$, the right derivative
of its enclosing-area infimum is
\begin{equation}
 a'_+(h)=\min_{F\in\mathcal F_*} a'_F(h),\qquad
 a'_F(h)=\frac12\int_{\partial^*F}
                       \operatorname{tr}_{T\partial^*F}h\,\dd A_g.
 \label{var:envelope-derivative}
\end{equation}
Here and below $\partial^*F$ is used for integration, while
``frontier'' means the support of its perimeter measure.
\end{proposition}

\begin{proof}
Large coordinate spheres have positive mean curvature, uniformly for
the finite-dimensional metric paths considered below.  On the region
outside one such sphere, the outward unit normal to their foliation
has positive divergence.  Applying the divergence theorem to the
portion of a filled competitor outside that sphere shows that clipping
off that portion cannot increase perimeter, and strictly decreases it
if the discarded portion has positive volume.  The argument holds for
almost every cutting radius by finite-perimeter slicing.  Choose two
fixed radii, both outside the compact core.  Clipping at the smaller
radius gives a common compact truncation, strictly inside the larger
one, for minimization and for all sufficiently small parameters of
each path.  BV compactness and lower semicontinuity now give a
minimizer; see \cite{AmbrosioFuscoPallara2000}.

A smooth full cut bounds a filled competitor with exactly the same
area, including its contact portions.  Conversely, first push $O$ and
an arbitrary filled competitor slightly outward by a smooth ambient
diffeomorphism supported in a collar of $S$.  The pushed competitor
contains a neighborhood of the original $O$, and its perimeter tends
to the original perimeter.  Strict approximation of finite-perimeter
sets by smooth domains can therefore be carried out while retaining
that neighborhood.  It is supported inside a fixed compact
truncation.  Retain the exterior component containing the end, filling
any bounded exterior components of the smoothed set.  This only
removes boundary components.  Its entire smooth intrinsic boundary is
a full cut.  Passing through these two approximations proves equality
of the infima.  The same removal observation shows that a minimizing
frontier cannot leave a removable bounded complementary open
component: a nonempty such component has positive perimeter, whose
removal would improve the minimum.  The usual indecomposable-component
decomposition gives the same conclusion for finite-perimeter
representatives \cite{AmbrosioFuscoPallara2000}.

Here is the local estimate needed at the obstacle.  For a local
competitor $G$, replace it by $G\cup O$.  Submodularity gives
\begin{equation}
 P(G\cup O)\le P(G)+P(O)-P(G\cap O).
 \label{var:obstacle-submodularity}
\end{equation}
Extend the smooth outward normal of $O$ to a smooth vector field of
norm at most one near its boundary.  Its divergence is bounded.
The calibration inequality for this field gives
$P(O)-P(G\cap O)\le C|O\setminus G|$ for changes in a sufficiently
small coordinate neighborhood.  Since $F$ minimizes among sets
containing $O$, it follows that
\begin{equation}
 P(F;B)\le P(G;B)+C|F\mathbin{\triangle}G|
 \label{var:quasiminimizing}
\end{equation}
whenever $F\mathbin{\triangle}G$ is compactly contained in the small
ball $B$.  This is an unconstrained perimeter almost-minimizing
estimate.  In coordinates the Riemannian area integrand is the
ellipsoidal integrand associated with
$\mathcal A=(\det g)g^{-1}$, which is smooth and uniformly elliptic.
The error in \eqref{var:quasiminimizing} is $O(r^n)$ and therefore is
bounded by $C_\gamma r^{n-1+\gamma}$ for any fixed
$0<\gamma<1$ on sufficiently small balls.  The almost-minimizer results in
\cite[Propositions~4.5--4.7 and Theorems~1.1, 9.5,
10.2--10.3]{Simmons2022} apply.  They give density and perimeter
bounds, local compactness, regularity off a set of dimension at
most $n-8$, and convergence as single $C^1$ graphs near free
regular limit points.  On free patches, the uniform local
$C^{1,\gamma/4}$ graph bounds and the smooth minimal-graph
equation upgrade this to smooth convergence on smaller patches
by interior elliptic estimates.  Tangent sets are
perimeter-minimizing cones.  The minimal-boundary formulation and
dimension reduction are also described in
\cite[Chapter 7, Section 5]{SimonGMT2014}.

At a contact point a tangent filled cone contains the tangent
half-space of $O$.  The corresponding minimizing hypercone is
therefore supported on one side of its boundary plane, and is that
plane.  One way to see the last assertion is to integrate the signed
coordinate-height equation on its stationary spherical link; the
height has one sign and its spherical Laplacian is a negative
constant times itself.  The singular set can be removed in this
integration by its dimension and the perimeter bounds.  Equivalently
one can use the half-space cone theorem in
\cite[Chapter~7, Section~4, Theorem~4.5]{SimonGMT2014}.  The almost-minimizer
regularity theorem now gives a $C^{1,\beta}$ graph at contact
for some $\beta>0$; one may take $\beta=\gamma/4$.
Estimate~\eqref{var:quasiminimizing}, differentiated along smooth
flows, bounds its weak mean curvature.  Difference quotients in the
uniformly elliptic prescribed-mean-curvature graph equation give
$W^{2,2}$ regularity.  Written in nondivergence form, that equation
has H\"older coefficients and bounded right-hand side, so local
$W^{2,p}$ estimates give all finite $p$; see
\cite{GilbargTrudinger2001}.

If two minimizing frontiers meet at a free regular point of one of
them, their union and intersection also minimize by
\eqref{var:obstacle-submodularity}, with equality in the sum of the
perimeters.  They are ordered relative to the given regular sheet.
Their tangent cones at contact are one-sided relative to its tangent
plane, and hence are planes.  Regularity and the strong comparison
principle for the minimal graph equation imply local coincidence.
Strict compactness and regular convergence then give the asserted
uniform graph neighborhoods: otherwise a sequence of failing
frontiers would converge strictly to a minimizer through the point,
where multiplicity-one regular convergence gives precisely the
missing neighborhood.  The difference of two distinct ordered
graphs solves a linear uniformly elliptic equation.  Positive-solution
Harnack and interior derivative estimates imply, on smaller patches,
that its gradient is bounded by a uniform constant times its
pointwise separation.

Finally let $F_t$ minimize for $g_t$.  Compactness, lower
semicontinuity and comparison with a fixed $g_0$ minimizer show that
every convergent subsequence at $t=0$ converges strictly in perimeter
to an element of $\mathcal F_*$.  Strict continuity for vector
measures then gives convergence of the tangent-plane area measures
\cite{AmbrosioFuscoPallara2000}.  The area Taylor expansion has a
uniform remainder on the resulting perimeter-bounded family.  Fixed
minimizers give the upper bound in
\eqref{var:envelope-derivative}; the varying minimizers give the lower
bound.  This also proves compactness of the minimizing family and
continuity of $F\mapsto a'_F(h)$ in its strict topology.
\end{proof}

\subsection{Strict directions and positive-measure separation}

An active ray is a pair $(x,v)$ with $|v|_g\le1$ and
$\mu_m(x)+J_m(x)(v)=0$.  Under a metric variation $h$, transport the
unit ball by the metric square root, so $v'=-h^\sharp v/2$.  Denote by
$\mathfrak C'$ the variation of $2(\mu_m+J_m(v))$ under that
transport.  Our test space consists of finite linear combinations of:
\begin{enumerate}[label=(\roman*)]
 \item arbitrary compact smooth variations $(h,p,\alpha',D_m')$,
 including support up to $S$, with $p=K'$ and both form variations
 closed;
 \item a metric end prototype $h=r^{2-n}\delta$, cut off toward the
 core;
 \item second-tensor end prototypes whose Euclidean trace reversal is
 \begin{equation}
  \Pi_{ij}=r^{-k}
   \bigl(a_i n_j^\delta+a_j n_i^\delta
                     -(a\cdot n^\delta)\delta_{ij}\bigr),
  \qquad n^\delta=x/r,
  \label{var:momentum-prototype}
 \end{equation}
 with arbitrary constant $a\in\R^n$ and a cutoff toward the core;
 \item the strict conformal direction $(h,p)=(2\phi g,\phi K)$,
 with both forms fixed.
\end{enumerate}
For the last direction choose $0<\beta<\min(1,q)$, put
$\rho_0=r^{-n-\beta}$ with a positive smooth extension of the radial
weight, and take
\begin{equation}
 -\Delta\phi=B\sqrt{|\dd\phi|^2+r^{-2k}}+\rho_0,
 \qquad \partial_\nu\phi=-1,
 \qquad \phi(\infty)=0,
 \label{var:strict-direction}
\end{equation}
where $B\ge |K|$ is smooth and has simple symbol bounds
$B=O_j(r^{-1-q})$.  The strictification construction in
Lemma~\ref{prep:strictification} applies to this decaying drift:
its radial barriers dominate the drift because $\beta<q$.  It gives
$\phi>0$, $\phi=O_2(r^{2-n})$ and finite gradient flux.  In
particular
\begin{equation}
 \frac{-\Delta\phi}{\rho_0}\longrightarrow1,
 \qquad
 \frac{\mathfrak C'_{\rm strict}}{\rho_0}>0
 \text{ on active rays},\qquad
 \frac{\mathfrak C'_{\rm strict}}{\rho_0}\longrightarrow2k,
 \qquad -\theta'_{\rm strict}=k.
 \label{var:strict-tail}
\end{equation}
Here the extra algebraic electric term is nonnegative: at a purely
electric active ray the conformal derivative of $\mathfrak C$ is
\begin{equation}
 -2k\Delta\phi+2kK(\nabla\phi,v)
                  +4(k-1)c_n^2\phi|E|^2.
 \label{var:conformal-active-derivative}
\end{equation}

The prototypes have zero flat linearized constraints on the end.
For example
$\Pi_{ij}=r^{-n}(a_ix_j+a_jx_i-(a\cdot x)\delta_{ij})$
satisfies $\partial_j\Pi_{ij}=0$ and
$\Pi_{ij}n^j=r^{-k}a_i$.  Its momentum derivative is $a_i/k$;
the scalar prototype has energy derivative $(n-2)/2$.
The background errors in the linearized constraints are
$O(r^{-n-q})$.  Thus the prototypes span the ADM flux derivatives
and their normalized active-constraint values tend to zero.

\begin{proposition}[Multiplier identity]
\label{var:multiplier-identity}
There are a probability measure $\pi$ on $\mathcal F_*$, a
nonnegative locally finite measure $\Lambda$ on the active rays,
and a nonnegative finite measure $\zeta$ on $S$, such that
\begin{equation}
 2k\omega\mathcal B'-c\int_{\mathcal F_*}a'_F(h)\,\dd\pi(F)
 =\int\mathfrak C'\,\dd\Lambda-\int_S\theta'_+\,\dd\zeta
 \label{var:separation}
\end{equation}
for every compact test and every end prototype above.  No
surjectivity of a constraint map is required.
\end{proposition}

\begin{proof}
Consider the linear image of the test space with components
\begin{equation}
 \left(\mathfrak C'/\rho_0,\ -\theta'_+,\
       \{c a'_F(h)-2k\omega\mathcal B'\}_{F\in\mathcal F_*}\right).
 \label{var:separation-map}
\end{equation}
The constraint test space is the one-point compactification of the
closed active-ray set, with one artificial infinity point added even
if that set is bounded.  At infinity assign $2k$ times the
coefficient of the strict direction.  The other two test spaces are
$S$ and the compact minimizing family.  The normalized functions
are continuous by \eqref{var:strict-tail} and
Proposition~\ref{var:area-envelope}.

This image does not meet the open cone of functions strictly
positive on the disjoint union of the three compact test spaces.
Indeed such a direction has positive strict-direction coefficient
because of the infinity test.  Apply that positive conformal
direction exponentially to the data after making the other
variations additively.  On compact subsets, strict first-order
positivity on the active rays implies feasibility for all
sufficiently small positive parameters: rays outside a fixed
neighborhood of the active set have a positive zeroth-order gap.
On the end the nonconformal linear errors are $O(r^{-n-q})$ and
are dominated by $\rho_0$.  After retaining the nonnegative
background through the conformal formula, the quadratic errors are
bounded by $Ct^2\rho_0$.  The strict first-order term therefore
gives $\mu_m-|J_m|\ge c_t\rho_0$, with $c_t>0$.
The forms remain closed; a newly created magnetic form has fixed
compact support.  Trapping is strict, the ADM vector remains
future timelike, and completeness, decay, integrability and
existence and differentiability of the fluxes are preserved.
These are precisely the hypotheses of
Corollary~\ref{flux:magnetic-extension}.  The area infimum stays
positive by uniform metric comparison on the compact core and
the controlled tail.  Equation~\eqref{var:envelope-derivative}
and strict positivity of the objective components would now
violate the numerical inequality for small positive parameter.

Hahn--Banach separation of a linear subspace from an open convex
cone gives a nonzero positive functional on the continuous
functions on this compact disjoint union.  The Riesz
representation theorem supplies the three nonnegative measures.
Its objective measure has positive mass.  Otherwise testing the
strict direction, uniformly positive on the two constraint test
spaces, would force both remaining measures to vanish as well.
Normalize the objective measure to have mass one.  On finite
active rays divide the constraint measure by $\rho_0$ to obtain
$\Lambda$.  This measure is finite over compact subsets,
including the boundary.  The possible infinity atom contributes
nothing on compact tests or prototypes.  Rearranging the
annihilation identity gives \eqref{var:separation}.
\end{proof}

\subsection{Normal-motion tests and the sheets selected by the multiplier}

Let $u$ and $X$ be the scalar and vector moments of $\Lambda$.
Initially these are measures, using the fixed Riemannian
volume-density convention; thus $u$ also denotes the scalar
distribution defined by its moment measure.  Positivity and the
definition of the active set imply
\begin{equation}
 u\ge |X|_g,\qquad u\mu_m+J_m(X)=0
 \label{var:moment-complementarity}
\end{equation}
as measure identities.  In particular, electromagnetic or
metric terms of the ray derivative require moments of order at
most one.

\begin{lemma}[Tangential trace on free sheets]
\label{var:trace-free-sheets}
For $\pi$-almost every minimizing frontier,
$\operatorname{tr}_{T\partial F}K=0$ on each of its free regular
sheets.  Every such frontier either equals $S$, or is wholly
inside $\operatorname{int}\Omega$ and satisfies
$H=\operatorname{tr}_{T\partial F}K=0$ on its regular part.
\end{lemma}

\begin{proof}
Fix a smooth lapse test $s$ compactly supported in the open
exterior, and choose a local Gaussian Lorentzian metric inducing
$g,K$ at time zero.  Its spatial metric and first normal jet are
fixed, while its second normal metric jets are free.  Choose a
two-form whose initial electric and spatial components are
$c_n E$ and zero.  Its first normal derivatives can be chosen so
that both Maxwell equations hold on the initial slice.  In
general dimension the tangential Maxwell constraints are
$\dd D_m=0$ and $\dd\alpha=0$; the remaining equations prescribe
the $n(n-1)/2$ magnetic and $n$ electric normal evolution
components.  The induced flux-form variations under normal
motion with velocity $sN$ are compactly supported and closed,
as follows also from Cartan's formula applied to these initial
Maxwell jets.

Let $\mathsf S$ denote the electromagnetic stress tensor with
the normalization in Definition~\ref{def:normalization}.
It is divergence-free at the initial slice.  For each
$\delta>0$, prescribe there the spatial components of
Einstein minus $\mathsf S$ to be
\begin{equation}
 T^{(\delta)}_{ij}
       =\frac{(J_m)_i(J_m)_j}{\mu_m+\delta}.
 \label{var:normal-jet-matter}
\end{equation}
Spatial trace reversal of the second normal metric jets is
invertible, so these are legitimate smooth jet prescriptions.
The normal and mixed components remain $\mu_m,J_m$ by
Gauss--Codazzi.  The resulting variation has $h=2sK$.

We check that $\mathfrak C'\to0$ uniformly on compact active
test sets.  The tensors in \eqref{var:normal-jet-matter}
converge spatially in $C^1$.  On $\{\mu_m>0\}$ this is
ordinary smooth convergence; near a zero of $\mu_m$, DEC gives
$|J_m|\le\mu_m$ and smoothness gives
$\dd\mu_m=\nabla J_m=0$ at that zero.  The bounds
\begin{equation}
 |T^{(\delta)}|\le\mu_m,
 \qquad
 |\nabla T^{(\delta)}|
       \le 2|\nabla J_m|+|\dd\mu_m|
 \label{var:matter-jet-bounds}
\end{equation}
make the convergence uniform across the zero set.
At a positive-density active ray, the limiting spacetime tensor
is $\gamma\otimes\gamma/\mu_m$, where
$\gamma(N)=\mu_m$ and $\gamma|_{T\Omega}=J_m$.
It is causal and annihilates the null vector $N+v$.
Parallel-transport that null vector along spatial curves.
Its contraction with the causal covector is nonnegative and
vanishes at the base point; differentiation at the minimum
therefore shows that all relevant spatial covariant
derivatives, with last argument $N+v$, vanish.  Bianchi and the
initial Maxwell equations then give the same limiting
vanishing for the normal derivative contracted with
$N,N+v$.  Changes of the first argument cost zero, and changes
of the second cost zero because it remains null and
$\gamma$ is proportional to its metric dual.  At a matter
zero, both the tensor and its spatial derivatives vanish by
\eqref{var:matter-jet-bounds}, and the divergence equation
again suffices.  These observations prove the asserted uniform
convergence of $\mathfrak C'$.

Apply \eqref{var:separation} to these tests and pass to the
limit.  They have no boundary or ADM variation, so
\begin{equation}
 \int_{\mathcal F_*}\int_{\partial^*F}
       s\operatorname{tr}_{T\partial^*F}K\,\dd A_g\,\dd\pi(F)=0.
 \label{var:weighted-trace-zero}
\end{equation}
At free regular points, Proposition~\ref{var:area-envelope}
provides uniform noncrossing graph patches, and frontiers
through a common point have the same tangent plane.
Consequently the tangential trace is a single-valued
continuous function on each such patch of the integrated
surface measure.  There is no cancellation between different
tangent planes in \eqref{var:weighted-trace-zero}.
A countable cover and Fubini imply the first conclusion,
initially almost everywhere and then everywhere on regular
sheets by smoothness.

At contact, the $W^{2,p}$ graph and the smooth obstacle have
the same second jets almost everywhere on their coincidence
set.  The frontier therefore satisfies
$H+\operatorname{tr}_{\rm tan}K=0$ almost everywhere throughout
the contact graph, including its free portion.  Uniform
ellipticity first improves its regularity and then bootstraps
it to a smooth solution of the expansion equation.  The
strong comparison principle gives coincidence with $S$ near
contact.  The set of coincidence is open and closed along
the connected $S$, so the whole $S$ is present.  Its area is
already the minimum $A$, leaving no area for another sheet.
If there is no contact, the compact frontier is separated
from $S$, and the stated equations hold on all its regular
points.
\end{proof}

\subsection{The moment equations and their initial regularity}

\begin{proposition}[Regularity of the measure moments]
\label{var:moment-regularity}
In the open exterior the moments are functions satisfying the first Killing initial data (KID) equation
\begin{equation}
 \operatorname{sym}\nabla X=-uK,
 \label{var:first-kid}
\end{equation}
and
\begin{equation}
 \operatorname{Hess}u
  =\mathcal L(u,X,\nabla X)
    +\frac c2\int_{\mathcal F_*}
       \nu_F^\flat\otimes\nu_F^\flat\,
                       \dd A_F\,\dd\pi(F).
 \label{var:hessian-measure}
\end{equation}
The equations are distributional, and surface measures in the
second equation are expressed relative to ambient volume.
The operator $\mathcal L$ has smooth linear coefficients.
It is specifically the lower-order operator from the metric
adjoint in \eqref{var:separation}, with the flux form fixed
and the unit-ball rays transported as specified above.
In particular, it is not an arbitrary operator with smooth
coefficients.  Locally $u$ is Lipschitz and $X$ is $C^{1,1}$.
The complementarity relations
\eqref{var:moment-complementarity} hold pointwise.
\end{proposition}

\begin{proof}
For a compactly supported interior variation $K'=p$, the
coefficient after integration by parts is
\begin{equation}
 2\bigl[u(\tau g-K)-\operatorname{sym}\nabla X
                                  +(\operatorname{div}X)g\bigr]:p.
 \label{var:second-tensor-adjoint}
\end{equation}
Its trace gives $\operatorname{div}X=-u\tau$, and its
trace-free part then gives \eqref{var:first-kid}.
For metric tests the scalar-curvature principal adjoint is
$\operatorname{Hess}u-(\Delta u)g$.
Connection variations in the momentum term give smooth
linear expressions in $u,X,\nabla X$; fixed-flux field terms
and the ray-transport term are algebraic in those moments.
The area term has coefficient
$-c(g-\nu_F^\flat\otimes\nu_F^\flat)/2$ on this side of
\eqref{var:separation}.  Inverting trace reversal gives
\eqref{var:hessian-measure}, with the displayed positive
normal-normal measure.

Differentiating \eqref{var:first-kid} and commuting
covariant derivatives expresses every component of
$\nabla^2X$ as a smooth linear combination of
$\operatorname{Rm}*X$, $K*\dd u$ and $(\nabla K)*u$.
Taking traces gives a strongly elliptic Laplace system for
$(u,X)$ with smooth lower-order coefficients.  Its source
measure has local growth $O(r^{n-1})$, uniformly in the
minimizer, by the local perimeter bound.  A localized
elliptic parametrix has order $-2$ and kernel bounds
$C|x-y|^{2-n-j}$ for $j=0,1$.  Such a potential of a measure
with this growth is locally $C^\gamma$ for every
$0<\gamma<1$: split the potential difference at twice the
point separation, use the growth bound nearby and the
first kernel derivative on successive outer annuli to get
$O(r(1+|\log r|))$.  Lower-order terms are included in the
parametrix, and the homogeneous remainder is smooth.
Thus $u,X$ are H\"older continuous.

The traced differentiated equation for $X$ now has the
form of the divergence of a H\"older field plus a
H\"older function; local elliptic estimates give
$X\in C^{1,\gamma}$, after decreasing $\gamma$ if needed.
The lower term in \eqref{var:hessian-measure} is therefore
continuous.  Its positive measure term implies a local
lower bound for the covariant Hessian in the distributional
sense.  Slice this inequality in smooth flow boxes whose
longitudinal curves are unit-speed geodesics.  Multiplying
by the smooth volume Jacobian gives the one-dimensional
distributional inequality $(u\circ\gamma)''\ge-C$ on
almost every such curve.  Continuity extends its integrated
convexity inequality to every curve in a smaller box.
This is local geodesic semiconvexity, which implies a local
Lipschitz bound by bounding one-sided slopes on short
geodesic segments with endpoints in a fixed larger ball.
The differentiated
equation for $X$ consequently has bounded right-hand
side componentwise, so $\nabla^2X$ is locally bounded and
$X\in C^{1,1}$.  The measure inequalities now reduce to
pointwise inequalities by continuity.
\end{proof}

For clarity about the field coefficients in this adjoint,
holding $\alpha=i_E\vol_g$ fixed gives
$E'=-(\operatorname{tr}h)E/2$ and hence
\begin{equation}
 \delta\bigl(2u c_n^2|E|^2\vol_g\bigr)
   =u c_n^2\bigl(2E_iE_j-|E|^2g_{ij}\bigr)h^{ij}\vol_g
   =-u\mathsf S_{ij}h^{ij}\vol_g.
 \label{var:electric-metric-adjoint}
\end{equation}
The ray transport contributes $-(J_m)_{(i}X_{j)}h^{ij}$.
The added variation of volume in the full constraint
pairing is zero by complementarity.  These identities,
together with the actual gravitational adjoint, will
control null focusing in the next section.  In
particular, the singular term in
\eqref{var:hessian-measure} cannot be discarded by
assuming regularity of a minimizing frontier.
